\section{Related Work \& Literature Landscape}
\label{sec:related_work}

The convergence of distributed ledger technology and healthcare informatics has inspired numerous system architectures over the past decade. In this section, we review prominent literature across medical blockchain systems, multi-chain distributed topologies, and privacy-preserving cryptographic storage models, concluding with a comparative taxonomic synthesis.

\subsection{Blockchain in Healthcare \& Medical Data Management}
Initial explorations into medical blockchains focused on decentralizing patient access controls and resolving the fragmentation of electronic medical records across institutional repositories. The seminal \textbf{MedRec} prototype developed by Ekblaw et al.~\cite{ekblaw2016medrec} demonstrated the utility of Ethereum smart contracts to manage authentication, identity resolution, and metadata pointers referencing off-chain provider databases. While MedRec successfully eliminated central identity authorities, it constrained its entire architecture to a single, monolithic blockchain instance where every participating medical provider was forced to synchronize all consortium state updates. Building upon this foundation, Dagher et al. introduced \textbf{Ancile}~\cite{dagher2018ancile}, an access control framework employing advanced proxy re-encryption on Ethereum to secure inter-organizational patient data sharing. Although Ancile achieved robust cryptographic privacy, it inherited the severe scalability ceilings of monolithic ledgers, lacking runtime execution isolation between high-throughput acute care centers and outpatient clinics. 

Subsequently, Zhang et al. pioneered the integration of standardized medical schemas with smart contracts in \textbf{FHIRChain}~\cite{zhang2018fhirchain}, demonstrating decentralized clinical research coordination and oncology trial cohort recruitment via HL7 FHIR payloads. However, FHIRChain relied on centralized cloud relayers to coordinate messaging and retained static consortium governance, leaving the network vulnerable to single-point administrative failure. Other notable contributions include \textbf{Healthcare Data Gateways} by Yue et al.~\cite{yue2016healthcare}, which provided mobile patient data ownership through multi-party computation, and \textbf{OmniPHR} by Roehrs et al.~\cite{roehrs2017omniphr}, which designed a distributed peer-to-peer personal health record index. Similarly, Xia et al. presented \textbf{MeDShare}~\cite{xia2017medshare} to monitor medical data provenance across cloud service providers, and Jiang et al. formulated \textbf{BloCHIE}~\cite{jiang2018blochie} to segregate personal healthcare telemetry from clinical diagnostic records across distinct blockchain tiers. Despite their respective advancements, these architectures either stored unencrypted hashes of immutable personal data on monolithic chains or relied on out-of-band administrative processes to govern network membership, failing to resolve the regulatory tension between data immutability and statutory erasure mandates.

[Note / Expansion Opportunity: Future analytical extensions could benchmark the smart contract gas overhead of proxy re-encryption schemes in Ancile against the native symmetric key lifecycle managed in our AES-256-GCM vault architecture.]

\subsection{Multi-Chain, Sidechain, and Sharding Architectures}
To overcome the throughput bottlenecks inherent in monolithic ledgers, distributed systems research has actively pursued sharded and heterogeneous multi-chain protocols. In the broader blockchain literature, the concept of a ``blockchain fork'' has historically carried a contentious and adversarial connotation, referring either to accidental soft forks arising from network latency or hard forks triggered by irreconcilable ideological and protocol splits, such as the Bitcoin/Bitcoin Cash and Ethereum/Ethereum Classic schisms~\cite{nakamoto2008bitcoin, wood2014ethereum, atzei2017survey}. In sharp contrast to these contentious divisions, cooperative forking has emerged as a promising strategy for application-specific domain specialization. Modern general-purpose multi-chain ecosystems, such as \textbf{Cosmos}~\cite{kwon2019cosmos} and \textbf{Polkadot}~\cite{wood2016polkadot}, facilitate horizontal scaling by interconnecting sovereign application-specific blockchains via standardized message passing protocols (e.g., IBC and XCMP). Furthermore, sharded systems such as \textbf{RapidChain}~\cite{zamani2018rapidchain} divide transaction processing and ledger state across randomized Byzantine committees. 

However, general-purpose multi-chain and sharding architectures exhibit critical limitations when transposed into healthcare environments. Neither Cosmos zones nor Polkadot parachains provide native genealogical mechanisms to inherit historical parent chain state at designated block heights ($\beta_{uv}$), forcing new clinical entities to either replicate entire global ledgers or start from an empty genesis block devoid of prior patient diagnostic baselines. Moreover, cross-chain communication protocols designed for token transfers (e.g., lock-and-mint bridge relays) are structurally ill-suited for complex clinical workflows, where an outpatient laboratory must cryptographically bind a diagnostic report to a preexisting physician order anchored on an independent hospital ledger.

[Note / Expansion Opportunity: A formal mathematical comparison could be developed proving that cross-chain communication complexity in an arborescence topology scales as $\mathcal{O}(|E|)$ compared to $\mathcal{O}(|V|^2)$ in fully meshed inter-blockchain networks.]

\subsection{Cryptographic Erasure \& Regulatory Compliance}
The collision between blockchain immutability and the European Union's General Data Protection Regulation (GDPR) has generated substantial debate in legal informatics and computer science~\cite{politou2018forgetting, finck2018blockchain}. Under GDPR Article 17, data subjects maintain the absolute legal right to obtain from the data controller the erasure of personal data concerning them without undue delay. Because the foundational premise of distributed ledger technology is the immutable permanence of historical blocks enforced by cryptographic hash linking, directly recording Protected Health Information (PHI) or identifiable clinical strings onto a distributed ledger creates an irreversible statutory violation. Several cryptographic strategies have been proposed to mitigate this tension, including zero-knowledge proofs (zk-SNARKs)~\cite{sasson2014zerocash}, ring signatures, and decentralized content-addressable storage networks such as the InterPlanetary File System (IPFS)~\cite{benet2014ipfs}. 

In particular, \textit{cryptographic erasure}—the deliberate destruction of the symmetric encryption keys required to decrypt off-chain ciphertexts—has emerged as the preeminent legal-technical compromise for blockchain architectures~\cite{politou2018forgetting, sharma2020security}. When the sole decryption key $\mathcal{K}$ for a ciphertext $c$ is expunged from all hardware security modules and key management systems, the residual ciphertext stored in off-chain vaults is rendered computationally indistinguishable from uniform random noise under the Indistinguishability under Chosen-Ciphertext Attack (IND-CCA2) security definition~\cite{bellare2000authenticated}. Furthermore, as established under the Random Oracle Model~\cite{bellare1993random}, the immutable on-chain cryptographic digest $h = \mathcal{H}(m)$ provides negligible advantage to any polynomial-time adversary seeking to invert the hash function or extract semantic clinical attributes. Consequently, the combination of authenticated symmetric encryption (AES-256-GCM)~\cite{dworkin2007recommendation} in off-chain vaults with immutable on-chain SHA-256 digest anchoring achieves both absolute auditability and verifiable compliance with GDPR Article 17 and HIPAA Security Rule \S~164.312~\cite{hipaa1996}.

[Note / Expansion Opportunity: Legal scholars can further evaluate whether multi-jurisdictional key escrow arrangements satisfy the rigorous definition of 'erasure' under recent Court of Justice of the European Union (CJEU) precedents.]

\subsection{State-of-the-Art Comparative Synthesis}
Table~\ref{tab:taxonomy} synthesizes our taxonomic evaluation, contrasting BlockchainForkTree against prevailing healthcare and multi-chain architectures across ten critical dimensions.

\begin{table*}[t]
\centering
\caption{Comprehensive Taxonomic Comparison of Healthcare Information Architectures}
\label{tab:taxonomy}
\footnotesize
\setlength{\tabcolsep}{2.5pt}
\begin{tabularx}{\textwidth}{p{3.1cm}*{6}{>{\centering\arraybackslash}X}}
\toprule
\textbf{Architectural Dimension} & \textbf{Centralized HIE} & \textbf{MedRec~\cite{ekblaw2016medrec}} & \textbf{FHIRChain~\cite{zhang2018fhirchain}} & \textbf{Hyperledger Fabric~\cite{androulaki2018hyperledger}} & \textbf{Cosmos IBC~\cite{kwon2019cosmos}} & \textbf{BlockchainForkTree (Ours)} \\
\midrule
Topology & Central Star & Monolithic Chain & Monolithic Chain & Peer Channels & Hub-and-Spoke & \textbf{Hierarchical Fork-Tree ($T$)} \\
Lineage Tracking ($\beta_{uv}$) & $\times$ None & $\times$ & $\times$ & $\times$ & $\times$ & \textbf{\checkmark (On-Chain Cryptographic)} \\
Governance Mechanism & Central Operator & Administrative & Central Admin & Fabric CA & Tendermint Voting & \textbf{\checkmark (On-Chain Consortium $\Phi(\Pi)$)} \\
Clinical Standard & Proprietary/HL7 & Custom JSON & HL7 FHIR & Custom Payloads & Interchain Tokens & \textbf{\checkmark (HL7 FHIR R4 Profiles)} \\
Storage Separation & Single Database & Off-chain Pointers & Off-chain Storage & Private Collections & State Tries & \textbf{\checkmark (AES-GCM + On-Chain SHA)} \\
GDPR Art.~17 Erasure & $\times$ Manual DB & $\times$ Hash persists & $\times$ Immutable CID & $\Delta$-Tombstones & $\times$ & \textbf{\checkmark (Provable Key Shredding)} \\
Fault Isolation & $\times$ Single POF & $\times$ Ledger Halt & $\times$ Ledger Halt & Partial (Channels) & \checkmark Per Zone & \textbf{\checkmark (Strict EVM Daemon Isolation)} \\
EHR Reconstruction & SQL Joins & Client Crawling & Single Contract & Cross-Channel SQL & Manual Relayer & \textbf{\checkmark ($\text{BFS-EHR}$ \& $\text{DFS-EHR}$ Proofs)} \\
RBAC Enforcement & Database ACL & Contract Modifiers & OAuth 2.0 Tokens & Fabric MSP & Tendermint Multisig & \textbf{\checkmark (Two-Tier Contract + Route Gating)} \\
Dynamic Fork Spawning & N/A & $\times$ & $\times$ & Manual CLI Config & Zone Governance & \textbf{\checkmark (Automated Multi-Party Consensus)} \\
\bottomrule
\end{tabularx}
\end{table*}

As detailed in Table~\ref{tab:taxonomy}, existing systems consistently falter when addressing the combination of genealogical lineage tracking, statutory erasability, and automated dynamic fork provisioning. Centralized HIEs completely lack decentralized fault tolerance and patient sovereignty. MedRec and FHIRChain suffer from monolithic throughput bottlenecks and lack formal mechanisms for GDPR key purging. Hyperledger Fabric provides channel-based partitioning but requires complex manual certificate authority orchestration that is vulnerable to administrative drift. Cosmos IBC excels in decentralized finance token passing but lacks native healthcare primitives, lineage block inheritance, and verifiable clinical order-fulfillment lifecycles. BlockchainForkTree is the first architecture to unite these disparate requirements into a coherent, mathematically grounded arborescence framework.

[Note / Expansion Opportunity: This taxonomy can be expanded into an empirical benchmark matrix comparing memory footprint and cold-start node deployment latency across Hyperledger Fabric versus BlockchainForkTree.]
